\subsection{Nonderogatory Matrices and Commutants}
\begin{definition} \label{def:nonderogatory-matrix} \hlt{Nonderogatory Matrix}\\
Let $n\in\mathbb{Z}_{>0}$, $A\in\operatorname{Mat}(n,n,\mathbb{C})$. Matrix $A$ \hlt{nonderogatory} if $(\forall\lambda\in\operatorname{spec}(A))$ its Jordan form contains exactly $1$ block with eigenvalue $\lambda$.\\
Equivalently, $(\forall\lambda\in\operatorname{spec}(A))\,\dim_{\mathbb{C}}E_{\lambda}(A)=1$ by Corollary \ref{cor:jordan-blocks-diagonalizability}(i).
\end{definition}

\begin{lemma} \label{lem:single-jordan-block-commutant-polynomials} \hlt{Commutant of One Jordan Block}\\
Let $m\in\mathbb{Z}_{>0}$, $\lambda\in\mathbb{C}$, $J\defeq J_m(\lambda)$, and $X\in\operatorname{Mat}(m,m,\mathbb{C})$. If $XJ=JX$, then $(\exists a_0,\ldots,a_{m-1}\in\mathbb{C})\,X=\sum\limits_{r=0}^{m-1}a_rJ^r$.
\begin{proof}
Set $N\defeq J_m(0)$; $J=\lambda I_m+N$, $XN=NX$ by $XJ=JX$. List $(\mathbf{e}_m,N\mathbf{e}_m,\ldots,N^{m-1}\mathbf{e}_m)$ is standard basis in reverse order.\\
Choose $b_0,\ldots,b_{m-1}\in\mathbb{C}$ satisfying $X\mathbf{e}_m=\sum\limits_{r=0}^{m-1}b_rN^r\mathbf{e}_m$ by uniqueness of basis coordinates.\\
Set $p(z)\defeq\sum\limits_{r=0}^{m-1}b_rz^r$. For $s\in\{0,\ldots,m-1\}$, $p(N)N^s=N^sp(N)$ by Proposition \ref{prop:operator-polynomial-products}(ii).\\
The equality $N^sX=XN^s$ holds by induction on $s$ from $NX=XN$, and $(p(N)-X)N^s\mathbf{e}_m=0$ by this equality and the choice of $p$.\\
Thus $p(N)=X$ by equality on a basis. Set $q(z)\defeq p(z-\lambda)$.\\
By binomial expansion, $(\exists a_0,\ldots,a_{m-1}\in\mathbb{C})\,q(z)=\sum\limits_{r=0}^{m-1}a_rz^r$. The equality $X=p(N)=q(J)=\sum\limits_{r=0}^{m-1}a_rJ^r$ holds by $J-\lambda I_m=N$.\\
\hlt{Structural Commentary}: Cyclic-basis equality propagates the image of $\mathbf{e}_m$ through every power of $N$.\\
\hlt{Key Insight}: Commutation determines the map from one cyclic vector.
\end{proof}
\end{lemma}

\begin{theorem} \label{thm:nonderogatory-commutant-polynomials} \hlt{Nonderogatory Commutants Are Polynomials}\\
Let $n\in\mathbb{Z}_{>0}$ and let $A\in\operatorname{Mat}(n,n,\mathbb{C})$ be nonderogatory. Then $\operatorname{Comm}(A)=\{p(A):p\in\mathbb{C}[z]\}$.
\begin{proof}
\begin{enumerate}[label=(\roman*)]
\setlength{\itemsep}{0pt}
\item \hlt{Polynomial Inclusion}: Fix $h\in\mathbb{C}[z]$ and regard $A$ as an operator on $\mathbb{C}^n$.\\
The equality $h(A)A=Ah(A)$ holds by Proposition \ref{prop:operator-polynomial-products}(ii) with the polynomials $h$ and $z$. Hence $\{p(A):p\in\mathbb{C}[z]\}\subseteq\operatorname{Comm}(A)$.
\item \hlt{Jordan Reduction}: By Theorem \ref{thm:complex-jordan-form}, choose $S\in\operatorname{GL}_n(\mathbb{C})$ and a Jordan matrix $J$ such that $A=SJS^{-1}$.\\
By nonderogatoriness, choose $d,m_1,\ldots,m_d\in\mathbb{Z}_{>0}$ and pairwise distinct $\lambda_1,\ldots,\lambda_d\in\mathbb{C}$.\\
For $i\in\{1,\ldots,d\}$, set $J_i\defeq J_{m_i}(\lambda_i)$. Then $J=\bigoplus\limits_{i=1}^{d}J_i$ and $\sum\limits_{i=1}^{d}m_i=n$.\\
Fix $B\in\operatorname{Comm}(A)$, set $X\defeq S^{-1}BS$, and partition $X=[X_{ij}]$ conformally with $J$.\\
Then $XJ=JX$ by multiplying $BA=AB$ on the left by $S^{-1}$ and on the right by $S$.\\
The equality $(\forall i\in\{1,\ldots,d\})\,\operatorname{spec}(J_i)=\{\lambda_i\}$ holds by Proposition \ref{prop:upper-triangular-characteristic-polynomial}(ii). Thus the spectra of $J_1,\ldots,J_d$ are pairwise disjoint.\\
Hence $(\forall i,j\in\{1,\ldots,d\})\,[i\neq j\Rightarrow X_{ij}=0]$ by Corollary \ref{cor:commuting-with-disjoint-blocks-has-zero-off-diagonal-blocks}.\\
Thus $X=\bigoplus\limits_{i=1}^{d}X_{ii}$. For $i\in\{1,\ldots,d\}$, choose $p_i\in\mathbb{C}[z]$ such that $X_{ii}=p_i(J_i)$ by Lemma \ref{lem:single-jordan-block-commutant-polynomials}.
\item \hlt{Congruence Projectors}: For $v,w,f\in\mathbb{C}[z]$ with $f\neq0$, write $v\equiv w\pmod f$ if $f\mid(v-w)$.\\
For $i\in\{1,\ldots,d\}$, set $f_i(z)\defeq(z-\lambda_i)^{m_i}$ and $F_i(z)\defeq\prod\limits_{j\neq i}f_j(z)$, where an empty product equals $1$.\\
Set $c_i\defeq F_i(\lambda_i)=\prod\limits_{j\neq i}(\lambda_i-\lambda_j)^{m_j}\neq0$ and $u_i(z)\defeq c_i^{-1}F_i(z)-1$.\\
The divisibility $z-\lambda_i\mid u_i(z)$ holds by $u_i(\lambda_i)=0$ and expansion in powers of $z-\lambda_i$. Set $e_i(z)\defeq c_i^{-1}F_i(z)\sum\limits_{r=0}^{m_i-1}(-u_i(z))^r$.\\
The identity $e_i=1-(-u_i)^{m_i}$ holds by $c_i^{-1}F_i=1+u_i$ and the finite geometric-sum identity. Hence $e_i\equiv1\pmod{f_i}$ by $z-\lambda_i\mid u_i$.\\
The relation $(\forall k\in\{1,\ldots,d\}\setminus\{i\})\,e_i\equiv0\pmod{f_k}$ holds by $f_k\mid F_i$.
\item \hlt{Block Interpolation}: Set $p\defeq\sum\limits_{i=1}^{d}p_ie_i$. Then $(\forall i\in\{1,\ldots,d\})\,p\equiv p_i\pmod{f_i}$ by item (iii).\\
For $i\in\{1,\ldots,d\}$, choose $q_i\in\mathbb{C}[z]$ satisfying $p-p_i=f_iq_i$. The identity $J_i-\lambda_iI_{m_i}=J_{m_i}(0)$ holds by Definition \ref{def:jordan-chain-block-and-matrix}(ii).\\
Hence $f_i(J_i)=J_{m_i}(0)^{m_i}=0$ by Lemma \ref{lem:nilpotent-jordan-block-power-rank}(i) and Definition \ref{def:column-space-matrix-rank-nullity}.\\
The equality $(p-p_i)(J_i)=f_i(J_i)q_i(J_i)=0$ holds by Proposition \ref{prop:operator-polynomial-products}(i). Thus $p(J_i)=p_i(J_i)=X_{ii}$.\\
The identity $p(J)=\bigoplus\limits_{i=1}^{d}p(J_i)$ holds by Proposition \ref{prop:direct-sum-matrix-algebra}(iv) and induction on $d$. Hence $p(J)=\bigoplus\limits_{i=1}^{d}X_{ii}=X$ by item (ii).
\item \hlt{Conjugation Back}: The identity $(\forall r\in\mathbb{Z}_{\geq0})\,A^r=SJ^rS^{-1}$ holds by induction on $r$ from $A=SJS^{-1}$.\\
Thus $p(A)=Sp(J)S^{-1}=SXS^{-1}=B$ by polynomial evaluation.\\
The claimed equality follows by the resulting reverse inclusion and item (i).
\end{enumerate}
\hlt{Structural Commentary}: Block separation isolates local polynomials; congruence projectors merge them before conjugation.\\
\hlt{Key Insight}: Truncated geometric sums realize every block congruence in one polynomial.
\end{proof}
\end{theorem}

Figure \ref{fig:nonderogatory-commutant-polynomial} tracks local commutant blocks into one polynomial satisfying every Jordan congruence.

\begin{illustrationfigure}[H]
  \centering
  \begin{tikzpicture}[
  font=\small,
  block/.style={draw=blue!60!black, fill=blue!6, rounded corners=2pt, align=center, inner sep=5pt, minimum width=2.00cm},
  poly/.style={draw=red!60!black, fill=red!6, rounded corners=2pt, align=center, inner sep=5pt, minimum width=2.33cm},
  arr/.style={-{Stealth[length=2mm]}, thick, gray!75}
]
  \node[block] (j1) at (0,0) {$J_{m_1}(\lambda_1)$};
  \node[block] (j2) at (2.55,0) {$J_{m_2}(\lambda_2)$};
  \node[block] (jd) at (5.10,0) {$\cdots$};
  \node[block] (jk) at (7.65,0) {$J_{m_d}(\lambda_d)$};

  \node[poly] (p1) at (0,-1.30) {$p_1(J_{m_1}(\lambda_1))$};
  \node[poly] (p2) at (2.55,-1.30) {$p_2(J_{m_2}(\lambda_2))$};
  \node[poly] (pd) at (5.10,-1.30) {$\cdots$};
  \node[poly] (pk) at (7.65,-1.30) {$p_d(J_{m_d}(\lambda_d))$};

  \node[poly, minimum width=3.52cm] (p) at (3.825,-2.52) {$p\equiv p_i\pmod{(z-\lambda_i)^{m_i}}$};

  \draw[arr] (j1) -- (p1);
  \draw[arr] (j2) -- (p2);
  \draw[arr] (jk) -- (pk);
  \draw[arr] (p1) -- (p);
  \draw[arr] (p2) -- (p);
  \draw[arr] (pk) -- (p);
\end{tikzpicture}

  \caption{Blockwise polynomial data merge into one polynomial modulo the Jordan-block powers.}
  \label{fig:nonderogatory-commutant-polynomial}
\end{illustrationfigure}

\begin{corollary} \label{cor:commutant-distinct-eigenvalues-polynomials} \hlt{Simple Spectrum Commutants Are Polynomials}\\
Let $n\in\mathbb{Z}_{>0}$ and let $A\in\operatorname{Mat}(n,n,\mathbb{C})$ have $n$ distinct eigenvalues. Then:
\begin{enumerate}[label=(\roman*)]
\setlength{\itemsep}{0pt}
\item $\operatorname{Comm}(A)=\{p(A):p\in\mathbb{C}[z]\}$;
\item $\dim_{\mathbb{C}}\operatorname{Comm}(A)=n$.
\end{enumerate}
\begin{proof}
List the eigenvalues as $\lambda_1,\ldots,\lambda_n$ and choose corresponding eigenvectors $\mathbf{v}_1,\ldots,\mathbf{v}_n$.\\
The list $(\mathbf{v}_1,\ldots,\mathbf{v}_n)$ is linearly independent by Theorem \ref{thm:distinct-eigenvalues-independent-eigenvectors}. It is a basis of $\mathbb{C}^n$ by Proposition \ref{prop:full-dimension-basis-criteria}(ii). Set $S\defeq[\mathbf{v}_1\ \cdots\ \mathbf{v}_n]$.\\
Then $S\in\operatorname{Mat}(n,n,\mathbb{C})^{\times}$ by Corollary \ref{cor:basis-columns-form-invertible-matrix}. The identity $A=S\operatorname{diag}(\lambda_1,\ldots,\lambda_n)S^{-1}$ holds by the eigenvector relations.\\
This Jordan form has one block per eigenvalue, so $A$ is nonderogatory by Definition \ref{def:nonderogatory-matrix}.\\
The equality $(\forall i\in\{1,\ldots,n\})\,\dim_{\mathbb{C}}E_{\lambda_i}(A)=1$ holds by Corollary \ref{cor:jordan-blocks-diagonalizability}(i).
\begin{enumerate}[label=(\roman*)]
\setlength{\itemsep}{0pt}
\item The equality $\operatorname{Comm}(A)=\{p(A):p\in\mathbb{C}[z]\}$ holds by Theorem \ref{thm:nonderogatory-commutant-polynomials}.
\item The equality $\dim_{\mathbb{C}}\operatorname{Comm}(A)=\sum\limits_{i=1}^{n}1^2=n$ holds by Theorem \ref{thm:commutant-of-diagonalizable-matrix}(ii).
\end{enumerate}
\hlt{Structural Commentary}: A full eigenbasis supplies nonderogatoriness; its eigenline dimensions feed the commutant formula.\\
\hlt{Key Insight}: A simple spectrum contributes one commutant parameter per eigenline.
\end{proof}
\end{corollary}
